\section{The term structure}
\label{sec:fol-metatheory-02}

Fix a maximal consistent \(\Delta\subseteq\Fm(\Lang^H)\) with the witness
property.  Define a relation on \(\Term(\Lang^H)\) by
\[
 r\sim_\Delta t
 \quad\text{iff}\quad
 \Delta\derives_{\ND_{\FOL}}r=t.
\]

\begin{lemma}[Equality congruence]
\label{lem:equality-congruence}
The relation \(\sim_\Delta\) is an equivalence relation.  If
\(n\in\N\), \(r_i,t_i\in\Term(\Lang^H)\), and
\(r_i\sim_\Delta t_i\) for \(1\leq i\leq n\), then for every
\(n\)-ary \(f\in\mathcal F\),
\[
 f(r_1,\ldots,r_n)\sim_\Delta f(t_1,\ldots,t_n),
\]
and for every \(n\)-ary \(R\in\mathcal R\),
\[
 R(r_1,\ldots,r_n)\in\Delta
 \quad\text{iff}\quad
 R(t_1,\ldots,t_n)\in\Delta.
\]
\end{lemma}
\begin{proof}
Reflexivity is \(=I\).  Symmetry was derived in
\cref{ch:fol-deduction}; transitivity follows by applying \(=E\) to the
formula \(r=x\).  Repeated applications of \(=E\), one argument at a time,
give compatibility with every function and relation symbol.  Maximal
consistent sets are deductively closed, so each derived replacement belongs
to \(\Delta\).  Symmetry supplies the reverse relational implication.
\end{proof}

\begin{definition}[Term structure]
The \emph{term structure} \(\Meta_\Delta\) has domain
\[
 \dom\Meta_\Delta
 =\Term(\Lang^H)/{\sim_\Delta}.
\]
Write \([t]\) for the class of \(t\).  Interpret symbols by
\[
 f^{\Meta_\Delta}([t_1],\ldots,[t_n])
 =[f(t_1,\ldots,t_n)]
\]
and
\[
 R^{\Meta_\Delta}([t_1],\ldots,[t_n])
 \quad\text{iff}\quad
 R(t_1,\ldots,t_n)\in\Delta.
\]
Let \(s_0\) be the variable assignment \(s_0(x)=[x]\).
\end{definition}

\Cref{lem:equality-congruence} proves that these interpretations are
independent of the chosen representatives.  The domain is nonempty
because \(\Lang^H\) has constant symbols.

\begin{lemma}[Term evaluation in the term structure]
\label{lem:term-model-evaluation}
For every term \(t\),
\[
 \sem{t}^{\Meta_\Delta}_{s_0}=[t].
\]
\end{lemma}
\begin{proof}
Induct on \(t\).  For a variable this is the definition of \(s_0\); the
function case is the definition of \(f^{\Meta_\Delta}\) applied to the
induction hypotheses.
\end{proof}

\begin{lemma}[Quantifier membership]
\label{lem:quantifier-membership}
For every formula \(\varphi\) and variable \(x\):
\begin{enumerate}
  \item \(\exists x\varphi\in\Delta\) iff
        \(\varphi[t/x]\in\Delta\) for some term \(t\); and
  \item \(\forall x\varphi\in\Delta\) iff
        \(\varphi[t/x]\in\Delta\) for every term \(t\).
\end{enumerate}
\end{lemma}
\begin{proof}
If an existential formula belongs to \(\Delta\), the witness property supplies
a constant instance.  Conversely, \(\exists I\) derives the existential from
any term instance, and deductive closure places it in \(\Delta\).

Universal elimination proves the forward implication in 2.  Conversely,
suppose \(\forall x\varphi\notin\Delta\).  Maximal consistency gives
\(\neg\forall x\varphi\in\Delta\); quantifier duality and deductive closure
give \(\exists x\neg\varphi\in\Delta\).  By the witness property,
\(\neg\varphi[c/x]\in\Delta\) for some constant \(c\), so consistency gives
\(\varphi[c/x]\notin\Delta\).  Therefore not every term instance belongs to
\(\Delta\).
\end{proof}

\begin{lemma}[First-order truth lemma]
\label{lem:fol-truth}
For every \(\Lang^H\)-formula \(\varphi\),
\[
 \Meta_\Delta,s_0\entails\varphi
 \quad\text{iff}\quad
 \varphi\in\Delta.
\]
\end{lemma}
\begin{proof}
Proceed by strong induction on \(c(\varphi)\).  For equality, term evaluation gives
\[
 \Meta_\Delta,s_0\entails r=t
 \quad\text{iff}\quad [r]=[t]
 \quad\text{iff}\quad \Delta\derives r=t
 \quad\text{iff}\quad r=t\in\Delta,
\]
where the last equivalence is deductive closure and the definition of
\(\sim_\Delta\).  The relation case follows directly from the interpretation
of \(R^{\Meta_\Delta}\), and \(\bottom\notin\Delta\) by consistency.  The
Boolean cases use the propositional closure properties of maximal consistent
sets exactly as in \cref{lem:prop-truth}.

For the existential case, term-structure evaluation and the substitution
lemma (\cref{thm:fol-substitution}) give, for every term \(t\),
\[
 \Meta_\Delta,s_0[x\mapsto[t]]\entails\varphi
 \quad\text{iff}\quad
 \Meta_\Delta,s_0\entails\varphi[t/x].
\]
By \cref{lem:substitution-complexity}, every instance
\(\varphi[t/x]\) has complexity \(c(\varphi)\), which is smaller than the
complexity of \(\exists x\varphi\).  The strong induction hypothesis and
\cref{lem:quantifier-membership}(1) therefore show that
some class \([t]\) satisfies the matrix iff \(\exists x\varphi\in\Delta\).
The universal case uses ``every'' and
\cref{lem:quantifier-membership}(2).  Every element of the term structure has the
form \([t]\), so all semantic quantifier cases are covered.
\end{proof}

